\cotefile{77}
\scanpage{129}

% L1
donc si $\unsure{Q'|q'}$ est lisse, $\delta'(q)$ définit un iso

% L2
\[
 L^{{\otimes}n}\overset{\sim}{\longleftarrow}(\det E)^{\otimes 2}.\tag*{(1.5)}
\]

% L3
Pour la suite on suppose que \unsure{$Q$ est|c'est|$q$ est} lisse

% L4
\uline{2} Groupe orthogonal\inserted{ et spécial-orthogonal}, cas $n$ impair $=2\nu+1$. // On a

% L5–6
\[
 \left\{\begin{aligned}
 q_n(x_1,y_1,\mathbin{-}x_\nu,y_\nu,z)&=\sum x_i y_i+z^2\\
 \delta'(q_n)&=(-1)^\nu \qquad \delta'(-q_n)=-\delta'(q_n)=(-1)^{\nu+1}
 \end{aligned}\right.
\]

% L6, left
On pose

% L7
\[
 O(q)=\underline{\mathrm{Aut}}(q)
\]

% L8–9
\diagram{\begin{tikzcd}[ampersand replacement=\&,column sep=small,row sep=small]
 {SO(q)=\underline{\mathrm{Aut}}(q)\cap\mathrm{Sl}(E)=\mathrm{Ker}\bigl(}
 \& {O(q)}\arrow[r,"\det"]\arrow[dr,bend right=20] \& {\mathbb G_m\bigl)}\\
 \& \& \mu_2\arrow[u,hook]
\end{tikzcd}}

% L10
\uline{NB} Sans condition sur parité de $n$, on a

% L11, left margin
\margin{\unsure{remonté|remonte|remarque}}

% L11
\[
 \mathbb G_m\cap O(q)=\mu_2\subset O(q)\qquad \mu_2=\underline{\mathrm{Centre}}\,O(q)\tag*{(2.1)}
\]

% L12
Dans le cas actuel $n$ impair, on a

% L13
\[
 \mu_2\cap SO(q)=\{1\}\,.\tag*{(2.2)}
\]

% L14–17, left rule
\uline{Thm} $O(q)$ est le produit direct des \uline{ss-groupes} $\mu_2$ et $SO(q)$, $SO(q)$ est lisse sur $S$ à fibres conn. (et en fait \deleted{simple si} réductif \deleted{si \unsure{$q$}} et $\unsure{=}\ \frac{1}{2}$ simple si $n\geq3$\inserted{ (i.e. $\nu\geq1$)}, simple si $n\geq7$\inserted{ (i.e. $\nu\geq\unsure{3}$)}).

% L18–20
Structure spéciale orthogonale sur $E$ \inserted{d'indice $\varepsilon$}, \deleted{\unsure{couple}} où $\varepsilon$ est une section de $\mathcal O_S^*$ (en pratique, $\varepsilon=\pm1$) : c'est une section-\inserted{base} $\omega$ de $\underline{\det}\,E$, telle que

% L21
\[
 \boxed{\underline{\omega}^{\otimes-2}=\varepsilon\,\delta'(q)}\tag*{(2.3)}
\]

% L21–24, right margin
\margin{Pour $n$ impair, \deleted{$q\mapsto\unsure{-q|q}$}\inserted{$(q,\omega)\mapsto(-q,\omega)$} \deleted{\unsure{fait passer}} bijection de l'ens des structures spéciales orth. pour $\varepsilon$, vers celles pour $-\varepsilon$ (donc pas g$^{\mathrm{de}}$ diff\ldots)}

% L22
On a évidemment

% L23
\[
 \mathrm{Aut}(E,q,\omega)=SO(q)\tag*{(2.4)}
\]

% L24–25
\[
 \left\{\begin{aligned}
 \text{Si}\qquad\varepsilon&=(-1)^\nu,\quad\text{alors }q_n\text{ est un prototype}\\
 \text{Si}\qquad\varepsilon&=(-1)^{\nu+1},\quad\text{---}\quad -q_n\quad\text{----------------}
 \end{aligned}\right.
\]

% lower left
% archive sheet number: 128
\scanpage{130}

% L1–2
$\lambda\mapsto\lambda^2$ sur le centre $\mathbb G_m$). Ici \deleted{pour} $n$ impair on trouve encore

% L3, slanting left margin
% slanted
\margin{commencer par ceci avant de prendre $O(q)$}

% L3–4
\uline{Prop} $O'(q)\overset{\sim}{\longleftarrow}\underbrace{SO(q)}_{\substack{\parallel\,(\mathrm{def}^{\mathrm n})\\O'(q)\cap\mathrm{Sl}(E)}}\times\mathbb G_m$.

% L5
qui équivaut au ---

% L6–11
\uline{Cor} (Pour $\varepsilon=\pm1$ fixé) La catégorie des modules quasi-orthogonaux \inserted{$(E,L,q)$ de rg $n$ impair $n$} sur $S$ est équiv.lente à la cat. des \deleted{\unsure{cples}} syst. $(E_0,q_0,\omega_0,\Delta)$, où $(E_0,q_0,\omega_0)$ est un module \inserted{spéc.} orth pour $\varepsilon$, de rang $n$, $\Delta$ est un module inv. À un système on associe

% L7–11, sideways left margin
% margin note written sideways
\margin{Dégager Thm plus \unsure{rigoureux|rigureux}}

% L12
\[
 E=E_0\otimes\Delta,\quad\inserted{L=\Delta^{\otimes 2},}\quad
 \varepsilon q_0^{\deleted{\unsure{\otimes\Delta|\otimes\Delta^{\otimes2}}}}:E_0\otimes\Delta\longrightarrow\Delta^{\otimes 2}=L,
\]

% L12, right margin
\margin{$\left(\begin{aligned}
 &\text{donc }\det E\simeq\overbrace{\det E_0}^{\mathrm{bse}}\otimes\Delta^n\\
 &\simeq\deleted{\det E_0\otimes{}}\Delta^{2\nu}\otimes\Delta\\
 &\simeq\deleted{\det E_0\otimes{}}L^\nu\otimes\Delta
\end{aligned}\right)$}

% L13
Le foncteur quasi-inv. associe à $(E,L,q)$

% L14, left margin
\margin{$n=2\nu+1$}

% L14
\[
 \Delta=\inserted{L^{\otimes-\nu}}\otimes\det E,\quad
 E_0=E\otimes\Delta^{-1},\quad q_0\deleted{\unsure{=|:}}\ \varepsilon q\otimes\Delta^{-1}:E_0=E\otimes\Delta^{-1}\longrightarrow L\Delta^{-2}
\]

% L15
\[
 \text{or }L\Delta^{-2}\simeq L\otimes L^{2\nu}\otimes(\det E)^{-2}\simeq L^n\otimes(\det E)^{-2}\simeq\mathcal O_S
 \quad\text{par }(\unsure{1.3|4.3})\quad\deleted{\text{par }(1.3)\quad\mathcal O_S}
\]

% L16
$\omega_0\in\Gamma(\det E_0)$ déduit de

% L17–20
\[
\begin{aligned}
 \det E_0&\simeq\underbrace{\det E}_{\Delta}\otimes\Delta^{-n}
 \deleted{=\Delta\,\unsure{\Delta^n\Delta^{2\nu+1}}\simeq L^\nu\Delta\quad\text{or}}\\
 &\deleted{\simeq\Delta L^\nu\Delta^{-n}}\\
 &\deleted{\det E_0\simeq\det E\otimes\Delta E^{\unsure{-\nu}}}\\
 &\simeq\det E\,L^{n\nu}(\det E)^{-n}=L^{n\nu}(\det E)^{-2\nu}\\
 &=(L(\det E)^{-2})^\nu=\mathcal O_S^{\otimes\nu}\simeq\mathcal O_S\qquad\unsure{\text{cf Thm}|\text{etc.}}
\end{aligned}
\]

% lower left
% archive sheet number: 129
\scanpage{131}

% L1–8
\uline{Prop} (Cor du Thm) Soit $\varepsilon=\pm1$. La cat. des modules orth. lisses \inserted{$(E,q)$} sur $S$ de rg $n=2\nu+1$ est équivalente à la cat. pro\deleted{\unsure{d}}duit des couples $(E_0,\deleted{\unsure{L}}q_0;\omega_0,\Delta,\varphi)$, où $(E_0,q_0,\omega_0)$ est un module spécial quadratique pour $\varepsilon$, \deleted{\unsure{et}} $\Delta$ un mod. inv., $\varphi:L^{\otimes2}\overset{\sim}{\longleftarrow}\underline O_S$ (NB le couple $(L,\varphi)$ revient à la donnée de un torseur sous $\mu_2$). Le foncteur établissant l'équivalence est

% L9
\diagram{\begin{tikzcd}[ampersand replacement=\&,column sep=small]
 {(E_0,\ldots)}\arrow[r,mapsto] \& {\bigl(E_0\otimes\Delta,\ q_0\otimes L:}
 \& {E_0\otimes\Delta}\arrow[r]\arrow[rr,bend right=15] \& {\Delta^{\otimes2}\deleted{\unsure{\otimes}}}\arrow[r,"{\overset{\varphi^{-1}}{\sim}}"] \& {\underline O_S\bigr)}
\end{tikzcd}}

% L10
Le foncteur quasi-inverse est

% L11
\[
 (E,q)\longmapsto\bigl(\Delta=\det E,\ \varphi=\delta'(q)^{-1},\ E_0=E\otimes\Delta^{-1},
\]

% L12
\diagram{\begin{tikzcd}[ampersand replacement=\&,column sep=small]
 {q_0=\varepsilon q\otimes\Delta^{-1}:E_0}\arrow[r]\arrow[rr,bend right=15] \& \Delta^{-2}\arrow[r,"{{}^{t}\varphi}"] \& {\mathcal O_S,}
\end{tikzcd}}

% L13–14
\[
 \begin{aligned}
 &\omega_0\in\Gamma\det E_0\deleted{=\Gamma(\det E)\otimes L^{-n}\simeq\underline O_S}\\
 &\text{déduit de }\det E_0\simeq
 \underbrace{(\det \deleted{\unsure{E}}\inserted{E})}_{\deleted{\det E}\,\Delta}\otimes
 \underbrace{\Delta^{-n}}_{\Delta^{\unsure{4|-1|2\nu|-4}}}\simeq\mathcal O_S
 \end{aligned}
\]

% L15
\uline{NB} $\mathrm{Aut}\,V(q)\simeq O(q)/\mu_2\simeq SO(q)$, il est lisse conn\inserted{, réductif},

% L16
\deleted{\unsure{NB}}\quad\deleted{\unsure{et $\Longrightarrow$}}

% L17–19
\textcircled{3} Structure\unsure{s} de similitudes (forme quadratique \inserted{lisse} donnée loc. mod. scalaire inv, ou encore forme lisse à valeurs dans un $L$

% L20
\[
 q:E\longrightarrow L
\]

% L21–24
Soit $O'(q)$\inserted{$\subset\mathrm{Gl}(E)\times\mathrm{Gl}(L)=\mathrm{Gl}(E)\times\mathbb G_m$,} le groupe de ses autom, \deleted{il en fait}\deleted{\unsure{en fait}} formé des couples $(u,\lambda)$ tels que $q(ux)\overset{\mathrm{idt}}{=}\lambda q(x)$ \deleted{$SO(q)$} $O'(q)\unsure{\text{---}|\longrightarrow}\mathrm{Gl}(\unsure{n|u|E})$ est injectif, $O'(q)$ contient $\mathbb G_m$ comme centre (l'hom $O'(q)\mathbin{-}\mathbb G_m$ induit

% lower left
% archive sheet number: 130
\scanpage{132}

% L1
\textcircled{4} \uline{Cas $n$ pair} $=2\nu$

% L2
\[
 q_n=q_{2\nu}=\sum_1^\nu(x_i y_i)\qquad \delta'(q_n)=(-1)^\nu\tag*{(4.1)}
\]

% L3
\deleted{\unsure{(4|(4.}} Ici on a

% L4–6
\[
 \underset{\substack{\parallel\\\mathrm{Centre}\,O(q)}}{\mu_2}
 \subset\underbrace{\deleted{\unsure{S}}O(q)\cap\mathrm{Sl}(E)}_{
 \substack{\text{on ne l'appelle}\\\text{\uline{pas}}\ SO(q)\\\text{à cause de car. }2}}
 \tag*{(4.2)}
\]

% L7
\deleted{et \unsure{on a}}

% L8–9, left rule
\uline{Thm} Pour $n$ pair, $O(q)$ est lisse, si $SO(q)$ est sa composante neutre, on a

% L10–11
\[
 \begin{aligned}
 &O(q)/SO(q)\simeq(\mathbb Z/2)_S\\
 \text{et}\qquad &\mu_2=\mathrm{Centre}(O(q))\subset SO(q).
 \end{aligned}
\]

% L12–13
\uline{Cor} Si $X=V(q)$, le groupe $\underline{\mathrm{Aut}}(X)\simeq\underline O(q)/\mu_2$ est lisse, et

% L14–15
\[
 \left\{\begin{aligned}
 &\underline{\mathrm{Aut}}(X)^{\unsure{0|\circ}}\simeq SO(q)/\mu_2\\
 &\underline{\mathrm{Aut}}(X)/\underline{\mathrm{Aut}}(X)^0\simeq(\mathbb Z/2)_S
 \end{aligned}\right.
\]

% L16
\unsure{Indication|Indications} de dém. du Thm.

% L17–18
La chose essentielle est de définir l'hom \deleted{\unsure{surj}} épi

% L18
\[
 O(q)\longrightarrow(\mathbb Z/2)_S
\]

% L19–20
et pour ceci on va associer can. à $q$ un rev. étale de degré $2$ de $S$.

% lower left
% archive sheet number: 131
\scanpage{133}

% L1–4, left rule
\uline{Prop} Considérons le sous-schéma de $\underline{\mathrm{Grass}}_\nu(E)$ \deleted{fermé aux}\inserted{correspondant aux} ``sous-espaces isotropes maximaux'' [sur lesquels $q$ \deleted{et la forme bil. associée sont} \inserted{est} \unsure{nuls|nul} ; $q(V)=0$,] \deleted{$\varphi\mid V\times E\longrightarrow\unsure{0}_S$}

% L5–7, left rule
C'est un sous-schéma fermé, il est lisse, et \deleted{pour} tte fibre géom, le nb de composantes connexes \deleted{est} $2$.

% lower left
% archive sheet number: 132
\scanpage{134}

% L1–2, title
\uline{Fibrés en drapeaux associés à}\deleted{ux}\inserted{ une} \uline{formes quadratique\unsure{s} régulière.}

% L3–5
$(E,q)$, avec $E$ de rg $N=2n+1$ (cas impair) ou $2n$, $n\geq1$ \inserted{donc $N\geq2$.} Soit $H\subset[1,n]\cap\mathbb Z$, \inserted{donc $H=\{i_1,\ldots,i_\nu\}$, où $1\leq i_{\unsure{1|j}}<\ldots<i_\nu\leq n$} \deleted{On pose} $\nu=\operatorname{card}H$

% L6–8
\[
 \begin{aligned}
 \Delta_H(E,q)&=\text{ens. des drapeaux de type }H\text{ de }E\\
 &E_1\subset E_2\ldots\subset E_\nu\quad(\nu=\operatorname{card} H,\ H=\{\operatorname{rg} E_1,\ldots,\operatorname{rg} E_\nu\}\\
 &\text{tels que }\deleted{\text{les}}\ E_\nu\text{ soit tot.}^{\mathrm t}\text{ isotrope.}
 \end{aligned}
\]

% L9
\[
 \underline\Delta_H(E,q):\mathrm{Sch}_{/S}\longrightarrow\mathrm{Ens}\qquad S'\longmapsto\Delta_H(E_{S'},q_{S'})
\]

% L10–12
\uline{Théorème 1} a) $\underline\Delta_H(E,q)\hookrightarrow\underline{\mathrm{Drap}}_H(E)$ est \deleted{représentable par} \inserted{une} immersion fermée -- en particulier $\underline\Delta_H(E,q)$ est représentable, et projectif sur $S$.

% L13–17
b) $\underline\Delta_H(E,q)$ est lisse sur $S$. Il est à fibres géométriquement $0$-connexes, sauf pour \deleted{$i_\nu=n$,} $\operatorname{rg}E$ pair $=2n$, $i_\nu=n$, auquel cas \deleted{\unsure{il}} ses fibres géom. ont exactement deux composantes connexes.

% L13–22, slanting left margin
% slanted
\margin{Donner aussi la structure de $\underline\Delta_H\to\underline\Delta_K$ pour $K\subset H$\ldots\ (Si $K=\varnothing$, on trouve $\underline\Delta_H\to S$\ldots) $\underline\Delta_H\to\underline\Delta_K$ est lisse, projectif, ses fibres géom sont $0$-conn. sauf si \deleted{\unsure{deux}} $\operatorname{rg} E$ est pair, $i_\nu=\operatorname{Max} H=n=\frac12\operatorname{rg} E$ \unsure{$\notin K$}, auquel cas les fibres géom. ont $2$ comp. conn.}

% L18–19
\uline{Dém} Le fait que $\underline\Delta_H(E,q)\to\mathrm{Drap}_H(E)$ une imm. fermée est immédiat, \unsure{..} .

% L20–22
\deleted{\unsure{Car si} $(E,q)$ est un}\inserted{Pour la lissité de} $\underline\Delta_H(E,q)$, on utilise le critère fonctoriel de lissité -- on est ramené à prouver que si $S_1\subset S$ est défini par $\mathcal J$ Idéal qu. coh. de carré nul,

% lower left
% archive sheet number: 133
\scanpage{135}

% L1–6
$S$ affine, et si $(E_{10},\ldots,E_{\nu0})$ est un élément de $\Delta_H(E_0,q_0)$, \unsure{ceci|cela|celui-ci} se remonte \inserted{en un $(E_1,\ldots,E_\nu)\in\Delta_H(E,q)$.} Il suffit bien sûr (\deleted{\unsure{par|par la}}\inserted{utilisant} les propriétés de lissité de $\underline{\mathrm{Drap}}_H\,E\longrightarrow\underline{\mathrm{Grass}}_{i_\nu}(E)$) de prouver que $F_0=E_{\nu0}$ se remonte en un sous-Module $F=E_\nu$ \deleted{\unsure{loc.}} fact. direct et \uline{tot. isotrope}.

% L6–8, cancelled working
\deleted{On procède par récurrence sur $d=\operatorname{rg}F$. Si $d=1$ on commence à remonter \unsure{en $F_0$} \unsure{en} $F$}

% L9–13, cancelled working
\deleted{($\underline{\mathrm{Grass}}_{i_\nu}(E)$ lisse sur $S$) facteur direct, mais $q\mid F$\deleted{\unsure{$0$}} est à valeurs dans $\mathcal J$, \inserted{donc peut s'identifier à un él. de $\mathrm{Quad}(F_0,\mathcal J)$} pas néc. nul, \deleted{On a}\inserted{Or} une ``correction'' sur le choix de $F$ dépend de}

% L14, cancelled display
\[
 \deleted{\xi\in\operatorname{Hom}(F_0,E_0/F_0\otimes\mathcal J)\simeq
 \Gamma\bigl(\underline{\mathrm{Hom}}(F_0,E_0/F_0)\otimes\mathcal J\bigr)}
\]

% L15–18, cancelled working
\deleted{puis supposons que $F_0$ a une base $(e_{i0})_{1\leq i\leq d}$ ($d=i_\nu$). On \deleted{\unsure{montre}}\inserted{note} que chaque $e_{i0}$ peut se remonter en $e_i$ \inserted{$\in\Gamma E$ tel que $q(e_i)=0$,} (lissité de la quadrique projective définie par $q$), soit}

% L19–23, cancelled working
\deleted{\deleted{$F$ le module engendré par les} c'est la lissité de la quadrique projective. Si $d>1$, \unsure{OPS} $F_0=G_0+\mathcal O_S e_0$, $G_0$ de rang $d-1$, $e$ section de $F_0$. On remonte $G_0$ en $G\in\Delta_{\{d-1\}}(E,q)$, puis il reste à remonter $e_0$ en section $e$ de $E$}

% lower left
% archive sheet number: 134
\scanpage{136}

% L1, cancelled
\deleted{satisfaisant}

% L2, cancelled display
\[
 \deleted{q(e)=0,\qquad\varphi(e,G)=0}
\]

% L3–4, cancelled
\deleted{Commençons par remonter un $e$ satisfaisant la \uline{première} de ces relations. Alors}

% L5, cancelled display
\[
 \deleted{x\mapsto\varphi(e,x)\qquad G\longrightarrow\underline O_S\quad\text{est à val. dans }\mathcal J}
\]

% L6, cancelled
\deleted{donc définit}

% L7, cancelled display
\[
 \deleted{\psi_0:G_0\longrightarrow\mathcal J\qquad\psi_0\in\underline{\mathrm{Hom}}(G_0,\mathcal J)}
\]

% L8, cancelled
\deleted{Essayons de corriger $e$ par}

% L9, cancelled display
\[
 \deleted{e'=e+\xi,\qquad\xi\in\Gamma(\mathcal J E)}
\]

% L10, cancelled
\deleted{la condition $q(e')=0$ s'écrit, compte tenu que}

% L11, cancelled display
\[
 \deleted{q(e)=0,\qquad q(\xi)+\varphi(e,\xi)=0}
\]

% L12, cancelled
\deleted{\deleted{ou encore}\inserted{or $q(\xi)=0$}}

% L13
Considérons l'hom

% L14–15
\[
 \begin{aligned}
 \underline{\mathrm{Hom}}(F_0,E_0/F_0)&\deleted{\unsure{\xi|\otimes\mathcal J}}\longrightarrow\underline{\mathrm{Quad}}(F_0)\\
 u_0&\longmapsto\bigl(x_0\mapsto\varphi_{\deleted{\unsure{F}}\unsure{0}}(x_0,u_0x_0)\bigr)
 \end{aligned}
\]

% L16–21
[ \uline{NB} $\varphi(F_0,F_0)=0$, donc $\varphi\mid F_0\times E_0$ se factorise par $F_0\times E_0/F_0$ ], il suffit de prouver que il est \deleted{surjectif}\inserted{épi} -- il restera épi par tensorisation par $\underline{\mathcal J}$, et en prenant sections sur base affine\ldots\ Or c'est un composé

% L22
\[
 \underline{\mathrm{Hom}}(F_0,E_0/F_0)\xrightarrow{\text{épi}}
 \underline{\mathrm{Hom}}\bigl(F_0,
 \underbrace{\overbrace{E_0/F_0^\perp}^{\substack{F_0^\vee\\\shortparallel}}}_{\deleted{\unsure{\sim|\simeq}F_0^\vee}}\bigr)
 \underset{\text{épi}}{\overset{\sim}{\longrightarrow}}\underline{\mathrm{Quad}}(F_0).
\]

% L22, above and below quotient
\inserted{isom : \unsure{explicitez|expliciter|explicitement} ($q$ régulier !) cf plus bas} \deleted{\unsure{isom à expliciter}}

% L22, left margin
\margin{\uline{NB}}

% lower left
% archive sheet number: 135
\scanpage{137}

% L1–6
Il faut rajouter le

% L2–6, left rule
\uline{Lemme} 1 Soit $F\subset E$, \deleted{loc. fact.}\inserted{$F\in\Delta_{\{d\}}(E,q)$}, Alors $F^\perp$ est loc. fact. direct et $\supset F$, et l'accouplement entre $F$ et $E/F^\perp$ est une dualité (donc $F=F^{\perp\perp}$, la formation de $F^\perp$ commute au chgt de base\ldots)

% L7–12
Il faut regarder \deleted{le noyau du \inserted{\uline{$F^\perp$}}} le composé $E\longrightarrow E^\vee\longrightarrow F^\vee$ et montrer que c'est \deleted{surjectif}\inserted{épi} -- le noyau étant $F^\perp$, le résultat s'ensuit. Or on est ramené au cas où $S$ est spectre d'un corps alg clos, [le seul cas non évident est celui où $\varphi$ \unsure{dégénérée|dégénère}, i.e. car $\unsure{k|h}=2$, $\operatorname{rg} E$ impair].

% L13–16
Mais \deleted{alors, si un $E\unsure{\text{---}|\to} E^\vee$ est le} si $a\in F$ est orth. à l'image de $E$, il est orth. à $E$ pour $\varphi$, \deleted{donc} ce qui joint à $q(\unsure{x|a})=0$ implique $a=0$, cqfd.

% L17–20
Reste à vérifier les assertions de connexion. Par dévissage, on est ramené, pour le cas \deleted{\unsure{$n$ p.}} $\operatorname{rg}E$ impair, ou $\operatorname{rg}E$ pair, $i_\nu<n$, \deleted{par}\inserted{le}

% lower left
% archive sheet number: 136

% L21–22
\uline{Lemme} \unsure{2|?} Soient $0\leq d<d'\leq n$, \inserted{\uline{NB} on admet ici $d=0$} \deleted{et supposons $d'<n$ si $\operatorname{rg}E$ pair}. Soit $F\in\Delta_{\{d\}}(E,q)$, consi%
\scanpage{138}%
% L1–3, left rule
dérons le sous-schéma de $\underline\Delta_{\{d,d'\}}(E,q)$ \deleted{défini par les $(F,G)$} image inverse de la section $F$ par $\underline\Delta_{\{d,d'\}}(E,q)\unsure{\text{---}|\to}\underline\Delta_d(E,q)$.

% L4–5, left rule
Ce schéma est \inserted{(lisse projectif)} \deleted{\unsure{à fibres non vides, $0$-connexes|à fibres non vides connexes}}, \deleted{\unsure{ses}}\inserted{à fibres géom. $0$-conn.} sauf si $\operatorname{rg}E$ pair et $d'=n$, auquel cas elles ont $2$ comp. conn.

% L6–7, cancelled argument
\deleted{\uline{Dém.} \inserted{\unsure{Sauf pour le résultat dans le cas} $\operatorname{rg}E$ pair, $d'=n$} \uline{OPS} par dévissage $d'=d+1$ (car l'image \inserted{continue} d'un $0$-connexe est $0$-connexe\ldots).}

% L8–9
Considérons $F^\perp$, (cf lemme 1) et considérons, \deleted{sur $S$ supposé affine, un scindage}

% L10
\[
 \deleted{F^\perp\simeq F\oplus\unsure{E'|G'}}\qquad E'=F^\perp/F
\]

% L11
muni de $q'$ déduit de $q$ par passage au quotient

% L12–16
\deleted{soit} \uline{Lemme 3} \deleted{On a $\unsure{G^\perp}\quad E\simeq G\oplus G^\perp$} \deleted{La forme} $E'$ est loc libre \deleted{et $G^\perp$ loc.} de rg $N-2d$ (donc de $\unsure{=|\text{même}}$ parité que le rg de $E$) et $q'$\deleted{\unsure{$E'$}} est régulière. Les $G\in\Delta_{\{d'\}}(E,q)$ tels que $F\subset G$ correspondent \unsure{biunivoq$^{\mathrm{mt}}$|biunivoquement} aux $G'\subset\Delta_{\{d'-d\}}(E')$.

% L17–18
Ceci nous ramène à déterminer les composantes connexes d'un $\Delta_{\{d\}}(E)$.

% L19–21
Ici encore, pour le cas de la $0$-conn., on est ramené par dévissage au cas de $d=1$, donc \inserted{\unsure{$\mathrm{rg}\,E$|de $\mathrm{rg}\,E$}} au \deleted{\unsure{lem}}

% L22–23, left rule
\uline{Lemme} 4 $\Delta_{\{1\}}(E)$ est \inserted{à fibres géom.} connexe si $\operatorname{rg} E\geq3$, ses fibres géom ont $2$ él. si $\operatorname{rg} E=2$

% L24
\uline{Évident} \deleted{\unsure{les|car}} en regardant dans $\mathrm P(E)$\ldots.

% lower left
% archive sheet number: 137
\scanpage{139}

% L1, cancelled heading
\deleted{(Dém. plus bas)}

% L2, upper text
Reste le cas $N=2n$, \inserted{et} $H=\{n\}$ i.e. $\underline{\Delta}_{\{n\}}(E,q)$.

% L3–7, upper text
\deleted{\unsure{On}} \deleted{Soit} Supposons d'abord 2 inv. sur $S$. Soit $F\in\Delta_{\{n\}}(E,q)$, \deleted{\uline{\unsure{Proposition}}} donc $F=F^\perp$ (via lemme 1), \deleted{d'où} $E/F\xrightarrow{\sim}F^\vee$, d'où un isom. can.

% L8, display with lower brace
\[
\deleted{\omega_F\,\unsure{:}}\quad
\underbrace{\det F\otimes\det F^\vee}_{\substack{\shortparallel\\ \underline{O}_S}}\longrightarrow\det E
\]

% L9–12, middle text
qu'on interprète comme une \uline{section} de $\det E$, notée $\omega_F$. Choisissons un suppl$^{\mathrm t}$ \inserted{$G$} (sur $S$ supposé affine) de $F$ dans $E$, donc

% L13, middle display
\[
E\simeq F\oplus G,\qquad \varphi\mid F\times G,\ \text{une dualité}
\]

% L14–16, lower text
Soit $(e_1\mathbin{-}e_n)$ une base de $E$ $(\exists\ \mathrm{loc.}\ (\mathrm{Zar}))$, $e_1^*\mathbin{--}e_n^*$ la \deleted{\unsure{f}}base duale $G\simeq F^\vee$, définie par

% L17, lower display
\[
\varphi(e_i,e_j^*)\simeq\delta_{ij}
\]

% L18, lower text
on a alors

% L19, bottom display
\[
\omega_F=e_1\wedge\mathbin{--}\wedge e_n\wedge e_1^*\wedge\mathbin{-}\wedge e_n^*
\]

% bottom left, pencil
% archive sheet number: 138
\scanpage{140}

% L1–4, bracketed lemma
[ \uline{Lemme} 5 ($S$ affine.) Soit $F\in\Delta_{\{n\}}(E,q)$, où $N=\mathrm{rg}\,E$ \deleted{\unsure{$N$}} $=2n$. Alors $\exists$ supp. $G$ de $F$ tel que $q\mid G=0$. (La forme $q$ sur $E\simeq F\times G$ s'exprime donc par

% L5–6, display
\[
\begin{aligned}
q(x\oplus y)&=\varphi\deleted{(}\langle x,y\rangle &&(=\langle x,y\rangle\ \text{si}\\
&&&\text{on identifie }G\unsure{:|\text{ à }}F^\vee)
\end{aligned}
\]

% L7–8, proof, bracket continues
\uline{Dém} On choisit un suppl$^{\mathrm t}$ quelc, puis on le corrige, en utilisant que

% L9, display and end of bracket
\[
\mathrm{Hom}\bigl(\underbrace{E/F}_{\substack{\shortparallel\\ \deleted{F}G}},\underset{\substack{\shortparallel\\G^\vee}}{F}\bigr)
\longrightarrow\mathrm{Quad}(G)\quad\text{est surjectif}\ldots\ \bigr]
\]

% L10–12, middle text
Donc dans le raisonnement précédent, \unsure{ops|ppg|OPS} $q\mid G=0$, de sorte que le calcul du discriminant donne

% L13, display
\[
\delta(q)=(-1)^n(\omega_F)^2
\]

% L14, text
ou encore

% L15, boxed display
\[
\boxed{(\omega_F)^2=(-1)^n\delta(q)^{-1}}
\]

% L16, cancelled line
\deleted{Considérons d'autre part}

% L17, lower text
On définit par $F\longmapsto\omega_F$ un morphisme

% L18–21, display and annotations
\[
\begin{aligned}
\underline{\Delta}_{\{n\}}(E,q)&\longrightarrow
M((-1)^n\delta(q))\overset{\substack{\text{\unsure{déf|dit} des variétés}\\\text{d'Arf}\ldots}}{=}\underline{\mathrm{Arf}}(E,q)\\[-2pt]
&\hspace{3em}\shortparallel\\[-2pt]
&\quad\text{sous-schéma de }V(\underline{\det}\,E)\\
&\qquad\text{des racines carrées de}\\
&\qquad (-1)^n\delta(q)^{-1}
\end{aligned}
\]

% bottom left, pencil
% archive sheet number: 139
\scanpage{141}

% L1, introduction
Nous pouvons maintenant énoncer

% L2–4, ruled theorem
\uline{Thm} 2 Supposons $\mathrm{rg}\,E=N$ pair $=2n$ \inserted{et supposons 2 inv. sur $S$.} Alors le morphisme précédent est lisse projectif à fibres \inserted{géom} $0$-connexes.

% L5–6, proof
Lisse projectif est évident, il faut voir que les fibres \inserted{géom} sont $0$-connexes.

% L7–10, cancelled proof
\deleted{Elles sont non vides : il suffit \inserted{(sur $k$ alg. clos)} de trouver \deleted{$\omega_{\unsure{F}}$} $F'$ avec $\omega_{F'}=-\omega_F$. Il suffit pour ceci de \deleted{remplacer} prendre}

% L10–11, cancelled display
\[
\begin{aligned}
\deleted{F'}&\deleted{=\deleted{\sum\unsure{\mathcal O}}\mathcal O_S e_1^*+\sum_{2\leq i\leq n}\underline O_S e_i}\\
\deleted{G'}&\deleted{=\mathcal O_S e_1+\sum_{2\leq i\leq n}\underline O_S e_i^*}
\end{aligned}
\]

% L12–17, lower text
Le fait que la flèche est surjective provient du fait que $\underline{\mathrm{Aut}}(E,q)=O(q)$ opère transitivement sur $\underline{\mathrm{Arf}}(E,q)$ (une \unsure{réflexion|reflexion} orthogonale p.ex. \deleted{\unsure{à déduire que}} \inserted{agit} comme \deleted{\unsure{symétrie fondamentale}} \inserted{l'opération canonique} du rev. d'Arf), il reste à prouver que les

% L18–20, lower continuation
\inserted{schéma} \inserted{$\underline\Delta_{\{n\}}^\omega(E,q)$} $F$ tels que $\omega_F=\omega$ (fixé) \deleted{forment} \unsure{forment|forme} un schéma connexe (sur $k$ alg. clos). Or c'est l'image de

% L21, bottom continuation
$\underline\Delta_{\deleted{\unsure{\{H\}}}\unsure{H\setminus\{1\}|\{1\}|\{1,\ldots,n\}}}^{\omega}$, où $H=[1,n]$, donc on est

% bottom left, pencil
% archive sheet number: 140
\scanpage{142}

% L1–3, continuation
ramené à prouver que ce dernier est connexe. Or il est fibré sur $\underline\Delta_{\{n-1\}}$ (qui est connexe), il suffit donc de prouver

% L4–5, lemma
\uline{Lemme 6} $\underline\Delta_{\{n-1,n\}}^\omega\xrightarrow{\sim}\underline\Delta_{\{n-1\}}$ par la projection canonique, en d'autres termes

% L5–7, left edge
% slanted
\margin{\unsure{indépendt|indépend|indépendant} de tte question de car.}

% L6–8, statement
$\forall F'\in\Delta_{\{n-1\}}(E,q)$, $\exists$ ex.ct$^{\mathrm{mt}}$ un $F\in\Delta_{\{n\}}^\omega(E,q)$ tel que $F\subset F'$. [Ou encore \deleted{\unsure{L}}la flèche de $(\underline\Delta_{\deleted{\unsure{H}}\,\{n-1,n\}})\mid F$ dans $\underline{\mathrm{Arf}}\deleted{\unsure{(E)}}$ est un iso]

% L9–14, proof
\uline{Dém} Par lemme 3, on voit que l'image inverse de $F'$ dans $\underline\Delta_{\{n,n-1\}}$ est $\unsure{\simeq|\text{à}}\Delta_1(\underbrace{F'^\perp/F'}_{\mathrm{rg}\,2})$ donc est un rev. étale de degré 2 de $S$. Il suffit de voir que le morphisme \inserted{de celui-ci} dans $\underline{\mathrm{Arf}}$ commute aux \unsure{autom.|actions} can. de ces rev., ce qui est immédiat.

% L15–20, lower text
Une fois traité le cas de car $\neq2$, \inserted{pour déterminer le nb de comp. \deleted{\unsure{de}} connexes de $\underline\Delta_{\{n\}}$,} on obtient le cas général par un argument de spécialisation (fact. de Stein....). En $=$ temps, \unsure{cela|on} définit $\underline{\mathrm{Arf}}(E,q)$ dans le cas général, par factorisation de Stein de $\underline\Delta_{\{n\}}$, et on

% L21, bottom continuation
voit que cela marche pour tous les $\underset{\substack{H\ni n\,\deleted{\unsure{\{n\}}}}}{\underline\Delta_H}$ \unsure{ou|où|avec}

% bottom left, pencil
% archive sheet number: 141
\scanpage{143}

% L1, introduction
Récapitulons

% L2–6, theorem and hypotheses
\uline{Thm} $(E,q)$ forme quadratique régulière de $\mathrm{rg}\ E=N$, $N=2n\deleted{\unsure{+1}}$ ou $N=2n+1$ ($n\geq1$ donc $N\geq2$). Soient $H,H'\subset[1,n]$, $\underset{\deleted{\unsure{d}}}{d}=\max H$, $d'=\max H'$ (donc $0\leq d,d'\leq n$), considérons

% L7, numbered display
\[
\underline\Delta_H(E,q)\xrightarrow{\unsure{p|\rho}_{H',H}}\underline\Delta_{H'}(E,q)\tag*{(1)}
\]

% L8, clause a
a) Ce morphisme est projectif, \deleted{et} lisse, \deleted{\unsure{à}} \inserted{surjectif}.

% L9–10, consequence
\deleted{b)} En particulier (prenant $H'=\varnothing$) les $\underline\Delta_H(E,q)\to S$ sont projectifs, lisses, surjectifs.

% L11–14, clause b
b) Les fibres géométriques de ce morphisme sont \deleted{\unsure{0}} connexes, sauf dans le cas $N$ pair, $d'<d=n$, auquel cas les fibres géométriques \deleted{\unsure{sont|ont}}\unsure{ont|sont} deux comp conn.

% L15, clause c
c) Dans le \deleted{dernier cas} \inserted{cas $N=2n$}, considérons

% L16–17, display
\[
\begin{aligned}
\underline{\mathrm{Arf}}(E,q)={}&\text{schém. des comp. conn}\\
&\text{de }\underline\Delta_{\{n\}}(E,q)\ \text{sur }S
\end{aligned}
\]

% L18–19, lower text
qui est un rev. étale de \deleted{degré} 2 de $S$. Considérons \inserted{(si $H\supset\{n\}$)} le \deleted{\unsure{fonct.|facteur}} morphisme composé

% L20, lower display
\[
\underline\Delta_H\longrightarrow\underline\Delta_{\{n\}}\longrightarrow\underline{\mathrm{Arf}}
\]

% L21, bottom text
\deleted{C'est lui qui \unsure{est ceci|est ici|ratisfait}} Ce morphisme est

% bottom left, pencil
% archive sheet number: 142
